Thanks to data augmentation, the spectroscopically informed datasets are highly representative and ideally suited for estimating photometric redshifts using machine learning, thus negating the need for predefined galaxy SED templates. However, future research into template-driven and hybrid methods would be valuable, despite being outside the scope of this study. 

In this section, we first introduce the \texttt{FlexZBoost} photometric redshift estimation method, and describe the definitions different galaxy sub-samples based on the point estimates of photometric redshifts. Although we use this specification here to characterise our photometric‐redshift quality, the SOM‐based augmentation method and the resulting datasets are applicable to many other science cases. 

\subsection{Machine learning modelling}
\begin{figure*}
    \centering
    \begin{subfigure}[b]{0.48\linewidth}
        \centering
        \includegraphics[width=\linewidth]{Figures/Y1/SAMPLE.pdf}
        \caption{LSST Y1}
    \end{subfigure}
    \hspace{0.0\linewidth}
    \begin{subfigure}[b]{0.48\linewidth}
        \centering
        \includegraphics[width=\linewidth]{Figures/Y10/SAMPLE.pdf}
        \caption{LSST Y10}
    \end{subfigure}
    \caption{Colour maps depicting the galaxy distribution, with the \(i\)-band magnitude on the y-axis and photometric redshift estimates, \(z_{\mathrm{phot}}\), on the x-axis, for LSST Y1 (left) and Y10 (right) under the baseline experiment. Black dashed lines indicate the selection range \(0.05 \le z_{\mathrm{phot}} \le 2.95\) used to identify source galaxies, while black dot-dashed lines represent the redshift-dependent magnitude limits imposed on lens galaxies \citep{2021PhRvD.103d3503P}.}
    \label{fig:7}
\end{figure*}

Based on a detailed comparative analysis with established methods \citep{2020MNRAS.499.1587S}, \texttt{FlexZBoost} \citep{10.1214/17-EJS1302, 2020A&C....3000362D} has demonstrated superior performance in estimating both the conditional probability density function (PDF) of photometric redshifts and the corresponding point estimates from galaxy colours. We therefore employ \texttt{FlexZBoost} in our proof‐of‐concept study to validate the efficacy of our SOM‐based data augmentation, deferring evaluation of alternative photometric‐redshift estimators to future work.  

\texttt{FlexZBoost} is built on the \texttt{Flexcode} framework, which models the conditional density \(p(z \mid \mathcal{T})\) for a given galaxy training dataset \(\mathcal{T}\) by projecting it onto an orthogonal basis of functions \(O_{k}(z)\) \(p(z \mid \mathcal{T}) \equiv \sum_{k} q_{k}(\mathcal{T}) \, O_{k}(z)\). Broadband photometry is therefore used to infer the expansion coefficients \(q_{k} (\mathcal{T}) = \int p(z \mid \mathcal{T}) \, O_{k}(z) \, \mathrm{d}z = \bigl\langle O_{k}(z) \mid \mathcal{T}\bigr\rangle\), which represent the conditional expectations of the basis functions. By employing regression in a high‐dimensional feature space, \texttt{FlexZBoost} effectively captures complex, non‐linear mappings between galaxy colours and redshift.  

Utilising the \texttt{FlexZBoost} implementation available on the \texttt{RAIL} platform, we develop our photometric redshift model using the \texttt{Combination} dataset, which integrates "observed" spectroscopic-like \texttt{Degradation} data with "simulated" \texttt{Augmentation} samples. This dataset is divided into training and validation sets in an \(80\% \,{:}\, 20\%\) proportion to perform regression on the hyperparameters of \texttt{FlexZBoost}, including the bump optimisation and sharpness parameters. The \(i\)–band apparent magnitude serves as the reference baseline against which all multi-band colours are defined.

We approximate the conditional PDF using a Fourier basis of up to \(50\) orthogonal functions uniformly covering \(0.0 \le z \le 3.0\) in steps of \(0.01\). Bump optimisation parameters are sampled over \([0.0, 0.5]\) on a \(50‐\)point grid, and sharpness parameters over \([0.5, 2.5]\) with the same grid resolution. For regression, we employ \texttt{XGBoost} \citep{2016arXiv160302754C}, minimising the Mean Squared Error (MSE) with a maximum tree depth of \(8\) and a learning rate of \(0.1\) -- choices that balance model expressivity with computational efficiency. Following hyperparameter tuning, the model is retrained on the complete \texttt{Combination} dataset. This rigorous framework provides both flexibility and robustness, effectively capturing the intrinsic complexity of galaxy populations while mitigating the risk of overfitting.

After training the \texttt{FlexZBoost} models using \(501\) experiments from the \texttt{Combination} dataset, each optimised model is applied to its respective experiment within the photometric \texttt{Application} datasets. This process generates the conditional density \(p \left( z | \mathcal{S}_n, \mathcal{P}_n \right)\) for each galaxy, with \(\mathcal{S}_n\) and \(\mathcal{P}_n\) representing the spectroscopic \texttt{Combination} and photometric \texttt{Application} datasets, respectively, for the \(n\mathrm{th}\) experiment. These conditions take into account variations in photometric noise, spatial coverage, and spectroscopic selection functions. We smoothly interpolate the conditional density over the interval \([0.0, 3.0]\) with redshift increments of \(0.01\), determining the point estimates \(z_{\mathrm{phot}} \) as the redshift that maximises the conditional PDF for each galaxy. 

\subsection{Definition of source and lens samples}
Based on point estimates, we first identify galaxies therefore serving as the background source galaxies that are gravitationally lensed by intervening mass distributions in the Universe \citep{2021A&A...645A.104A, 2022PhRvD.105b3514A, 2022PhRvD.105b3515S, 2023A&A...679A.133L, 2023PhRvD.108l3519D, 2023PhRvD.108l3518L, 2023OJAp....6E..36D, 2025arXiv250319442S, 2025arXiv250319441W}, therefore targeting galaxies with reliable measurements on weak gravitational lensing signal. Initially, we estimate the shear measurement noise, \(\sigma_\gamma\), for individual galaxies using the expression provided in \citep{2002AJ....123..583B}, written as 
\begin{equation} \label{eqa:3}
    \sigma_\gamma = \frac{a}{\mu} \left[ 1 + \left( \frac{b}{\eta} \right)^c \right].
\end{equation} The coefficients in the equation are \(\left(a,\, b,\, c\right) = \left(1.58,\, 5.03,\, 0.39\right)\) as defined by \cite{2013MNRAS.434.2121C}. In this calculation, as we described in Section~\ref{sec:2}, the variable \(\mu\) represents the \(\mathrm{SNR}\) of coadded \(r-\) and \(i-\)band photometry, while \(\eta\) signifies the squared ratio of the galaxy sizes of galaxies comparing to those of PSFs. Subsequently, we implement a selection criterion of \(\sigma_\gamma < \sigma_\epsilon\), where \(\sigma_\epsilon\) represents the intrinsic shape noise, presumed to be \(0.26\). 

On top of the constraints on the shear measurement noise \(\sigma_\gamma\), we further restrict our sample to galaxies with photometric redshifts in the range \(0.05 < z_{\mathrm{phot}} < 2.95\). This excludes objects at the boundaries of the redshift PDF, where asymmetries can lead to significant errors, and minimises contamination from higher-redshift (\(z > 3.0\)) galaxies, which contribute little to the lensing signal but complicate redshift calibration in a real analysis.  

Additionally, we define another foreground galaxy subset, referred to as lens samples, selected via precise redshift estimations to trace high‐mass haloes for large‐scale structure studies\citep{2021A&A...646A.140H, 2022PhRvD.105b3520A, 2023A&A...675A.189D, 2023PhRvD.108l3521S, 2023PhRvD.108l3517M}. We adopt the redshift‐dependent magnitude cut of \citet{2021PhRvD.103d3503P}, \( i < 4 \,\cdot\, z_{\mathrm{phot}} + 18 \), valid for \(0.2 < z_{\mathrm{phot}} < 1.2\). For both LSST Y1 and Y10, we apply the same criterion in accordance with the stringent photo-\(z\) precision requirements of the LSST DESC SRD. Future work should explore the inclusion of fainter galaxies in the Y10 lens samples -- while increased number density may enhance statistical power, it could also degrade redshift accuracy -- necessitating further investigation to optimise selection strategies \citep{2023ApJ...950...49M}.  

Figure~\ref{fig:7} shows the distribution of galaxies in \(i\)-band magnitude versus photometric redshift \(z_{\mathrm{phot}}\). Black dashed lines denote the source sample selection boundaries, which remove objects near analysis limits to prevent bias in the conditional density estimates, while black dot–dashed lines indicate the lens sample selection, isolating brighter galaxies to secure high photo-\(z\) precision. The distribution becomes increasingly disordered at faint magnitudes, highlighting the challenge of moving from training on \texttt{Degradation} datasets derived from \texttt{OpenUniverse2024} to using \texttt{Augmentation} datasets drawn from \texttt{CosmoDC2}. In the Y1 case, the spectroscopic selection limit, especially for \(20  \lesssim i_\ast \lesssim 22 \), noticeably biases redshift estimates at \(z_{\mathrm{phot}} \gtrsim 1.0\). By contrast, the deeper spectroscopic coverage in Y10 mitigates this bias, underscoring the importance of consistent lens sample definitions across survey epochs.  